% ============ 32 Lefschetz 铅笔的上同调 (原书 pp.186-192) ============
\chapter{Lefschetz 铅笔的上同调}\label{sec:32}
\index{消失循环}
\index{Picard-Lefschetz 公式}
\index{驯基本群}
\index{惯性群}
\index{Kazhdan-Margulis 定理}
\index{算术亏格}

本节中，$k$ 是特征 $p$(可能为 0)的代数闭域.

\subsection*{局部常值层的预备}

设 $U$ 是完备曲线 $X$ 的开子集，$\Lambda$ 是交换环(例如 $\Lambda = \mathbb{Z}/\ell^{n}\mathbb{Z}$, $\Lambda = \mathbb{Z}_{\ell}$, $\Lambda = \mathbb{Q}_{\ell}$).我们已经看到，给出 $U$ 上 $\Lambda$-模的可构造局部常值层，就是给出带 $\pi_1(U)$ 连续作用的有限生成 $\Lambda$-模.

当我们考虑 $X$ 上 $\Lambda$-模的可构造层时，情形更复杂：可以几乎任意地改变层在 $X \setminus U$ 的点上的茎(见命题~\ref{prop:8.17}).然而，有一类特殊的这样的层，对它们 $F$ 由 $F|_{U}$ 决定，即典范映射 $F \to j_{*}j^{*}F$ 是同构的那些层.如常，$j$ 是包含 $U \hookrightarrow X$.

\begin{proposition}\label{prop:32.1}
设 $U$ 是代数闭域 $k$ 上完备非奇异曲线 $X$ 的开子集，$\Lambda = \mathbb{Z}/\ell^{n}\mathbb{Z}$，$\mathbb{Z}_{\ell}$ 或 $\mathbb{Q}_{\ell}$，其中 $\ell$ 与 $k$ 的特征互素.

(a) 设 $F$ 是 $X$ 上 $\Lambda$-模的层；典范映射 $F \to j_{*}j^{*}F$ 是同构，当且仅当对所有 $s \in X \setminus U$，逆特殊化映射 $F_s \to F_{\bar{\xi}}$ 单射，且像为 $F^{I_s}_{\bar{\xi}}$(这里 $\bar{\xi}$ 是一般几何点，$I_s \subset \pi_1(U, \bar{\xi})$ 是 $s$ 处的惯性群).

(b) 对 $U$ 上 $\mathbb{Q}_{\ell}$-向量空间的任何局部常值层 $F$，杯积配对
\[
H^{r}(X, j_{*}F) \times H^{2-r}(X, j_{*}\mathcal{L}F(1)) \to H^{2}(X, \mathbb{Q}_{\ell}(1)) \cong \mathbb{Q}_{\ell}
\]
非退化.
\end{proposition}

\begin{proof}
(a) 映射 $F \to j_{*}j^{*}F$ 对所有 $x \in U$，在茎上诱导同构.对 $s \in X \setminus U$，两个层的茎分别是 $F_s$ 与 $F^{I_s}_{\bar{\xi}}$，因此该映射在 $s$ 处诱导同构，当且仅当逆特殊化映射是同构 $F_s \to F^{I_s}_{\bar{\xi}}$.

(b) 这可以由通常的 Poincar\'e 对偶定理~\ref{thm:14.7} 通过局部计算推出.见 EC V 2.2b.
\end{proof}

\subsection*{驯基本群}

\textbf{局部情形} 设 $K = k((T))$，即 $k$ 上的 Laurent 级数域.则 $K$ 是完备赋值环 $k[[T]]$ 的分式域，因此可以像数论中那样研究它的扩张域.在一点上这个域更简单：由于剩余域代数闭，没有非分歧扩张.于是，对 $K$ 的有限扩张 $L$，分歧次数等于次数.特别地，$K$ 的驯分歧扩张恰好是次数与 $p$ 互素的扩张.驯分歧扩张恰好是 $K$ 的形如 $K[T^{\frac{1}{d}}]$($\gcd(p, d) = 1$)的 Kummer 扩张.

固定 $K$ 的代数闭包 $K^{\mathrm{al}}$，设 $K^{\mathrm{tame}}$ 是 $K^{\mathrm{al}}$ 中所含 $K$ 的所有驯扩张的复合.对每个不被 $p$ 整除的 $\gamma \in \operatorname{Gal}(K^{\mathrm{tame}}/K)$ 与 $d$，
\[
\gamma T^{\frac{1}{d}} = \zeta T^{\frac{1}{d}}, \qquad \text{某个 } \zeta \in \mu_d(k).
\]
于是，有良定义映射 $\operatorname{Gal}(K^{\mathrm{tame}}/K) \to \mu_d$.对 $d$ 取极限，得到同构
\[
t \colon \operatorname{Gal}(K^{\mathrm{tame}}/K) \to \prod_{\ell \neq p} \mathbb{Z}_{\ell}(1).
\]

\textbf{整体情形} 现在设 $K = k(T)$，即 $k$ 上符号 $T$ 的有理函数域.则 $K$ 是 Dedekind 整环 $k[T]$ 的分式域，可以像数域情形那样研究它的扩张.$k[T]$ 的素理想形如 $(T - a)$($a \in k$ 适当)，而 $k(T)$ 关于由 $(T - a)$ 定义的赋值的完备化是 $k((T - a))$.还有一个额外的素理想，即对应 $k[T^{-1}]$ 中素理想 $(T^{-1})$ 的``无穷远素理想''.设 $S$ 是 $k[T]$ 的有限素理想集.

固定 $K$ 的代数闭包 $K^{\mathrm{al}}$，设 $K^{\mathrm{tame}}$ 是 $K^{\mathrm{al}}$ 中所含所有这样的有限扩张 $L/K$ 的复合：$L$ 在 $k[T]$ 的所有不在 $S$ 中的素理想上非分歧，在 $s \in S$ 的所有素理想上驯分歧.设 $K_s$ 是 $K$ 在素点 $s \in S$ 处的完备化，$K^{\mathrm{tame}}_s$ 是 $K_s$ 的最大驯分歧扩张.$K$ 到 $K_s$ 的嵌入到 $K^{\mathrm{tame}}$ 的延拓的选取，即交换图
\[
\begin{tikzcd}
K^{\mathrm{tame}} \arrow[r] & K^{\mathrm{tame}}_s \\
K \arrow[u] \arrow[r] & K_s \arrow[u]
\end{tikzcd}
\]
的选取，确定单射同态 $\operatorname{Gal}(K^{\mathrm{tame}}_s/K_s) \to \operatorname{Gal}(K^{\mathrm{tame}}/K)$.它的像 $I_s$ 由 $s$ 唯一确定到共轭.对每个 $s \in S$，有 $\operatorname{Gal}(K^{\mathrm{tame}}/K)$ 的子群 $I_s$ 与满同态 $t \colon I_s \to \mathbb{Z}_{\ell}(1)$.子群 $I_s$ 不必生成 $\operatorname{Gal}(K^{\mathrm{tame}}/K)$，因为可能存在 $K$ 的在 $k[T]$ 的所有素理想上非分歧，但在无穷远素理想上野分歧的真扩张(见例~\ref{ex:3.2}).

\textbf{几何情形} 设 $U$ 是 $k$ 上射影直线 $\mathbb{P}^{1}$ 的开子集，$S$ 是 $U$ 在 $\mathbb{P}^{1}$ 中的补集.我们希望研究在 $U$ 的点上非分歧，在 $S$ 的点上驯分歧的有限映射 $V \to \mathbb{P}^{1}$.本质上这与最后一种情形相同，只是多了一个对应无穷远点的素理想.

固定 $K \overset{\mathrm{def}}{=} k(\mathbb{P}^{1})$ 的代数闭包 $K^{\mathrm{al}}$，设 $\bar{\xi} \to U$ 是对应的几何点.设 $\pi^{\mathrm{tame}}_1(U, \bar{\xi}) = \operatorname{Gal}(K^{\mathrm{tame}}/K)$，其中 $K^{\mathrm{tame}}$ 是 $K^{\mathrm{al}}$ 中在对应 $U$ 的点的所有素理想上非分歧，在对应 $S$ 的点的素理想上驯分歧的子域的复合.则 $\pi^{\mathrm{tame}}_1(U, \bar{\xi})$ 对每个 $s \in S$ 含子群 $I_s$，且现在由这些子群生成(见例~\ref{ex:3.2}).对每个 $s$，有典范满同态 $t \colon I_s \to \mathbb{Z}_{\ell}(1)$.

\subsection*{上同调}

我们希望研究上同调在 Lefschetz 铅笔中如何变化，但首先应该看曲线的情形.本小节中，基域 $k$ 代数闭.

\textbf{曲线的铅笔} 考虑紧合平坦映射 $\varphi \colon X \to \mathbb{P}^{1}$，其纤维是不可约曲线，除 $s$ 在有限集 $S$ 中外非奇异，且对 $s \in S$，$X_s$ 以单个结点为奇点.则
\[
\chi(X_t, \mathcal{O}_{X_t}) \overset{\mathrm{def}}{=} \dim_k H^{0}(X_t, \mathcal{O}_{X_t}) - \dim_k H^{1}(X_t, \mathcal{O}_{X_t})
\]
在族中恒定(Mumford, \emph{Abelian Varieties}, p.50)，且对所有 $t$，$H^{0}(X_t, \mathcal{O}_{X_t}) = k$(因为 $X_t$ 不可约).因此，$\dim_k H^{1}(X_t, \mathcal{O}_{X_t})$ 恒定，记为 $g$.按定义，$\dim_k H^{1}(X_t, \mathcal{O}_{X_t})$ 是 $X_t$ 的算术亏格.当 $X_t$ 非奇异时，它是通常的亏格；当 $X_t$ 有单个结点时，$X_t$ 的正规化的亏格是 $g - 1$(Serre, J.-P., \emph{Groupes Alg\'ebriques et Corps de Classes}, Hermann, 1959, IV.7，命题 3).

当 $X_t$ 非奇异时，$H^{1}(X_t, \mathbb{Q}_{\ell})$ 的维数为 $2g$(见命题~\ref{prop:14.2}).与非奇异情形一样，为计算 $H^{1}(X_s, \mathbb{Q}_{\ell})$($s \in S$)，首先必须计算 $X_s$ 的 Picard 簇.

设 $\varphi' \colon \widetilde{X}_s \to X_s$ 是 $X_s$ 的正规化.它是同构，只是 $\widetilde{X}_s$ 的两个点 $P_1$ 与 $P_2$ 映到 $X_s$ 的奇异点 $P$.映射 $f \mapsto f \circ \varphi'$ 把 $P$ 的开邻域 $U$ 上的函数等同于 $\varphi'^{-1}(U) \subset \widetilde{X}_s$ 上在 $P_1$ 与 $P_2$ 处取相同值的函数.因此，有 $(X_s)_{\mathrm{et}}$ 上层 的正合序列
\[
0 \to \mathbb{G}_{m, X_s} \xrightarrow{f \mapsto f \circ \varphi'} \varphi'_*(\mathbb{G}_{m, \widetilde{X}_s}) \xrightarrow{f \mapsto \frac{f(P_1)}{f(P_2)}} \mathbb{G}_{m, P} \to 0.
\]
它的上同调序列是
\[
0 \to \mathbb{G}_{m} \to \operatorname{Pic}(X_s) \to \operatorname{Pic}(\widetilde{X}_s) \to 0,
\]
从中可以提取正合序列
\[
0 \to \mathbb{G}_{m} \to \operatorname{Pic}^{0}(X_s) \to \operatorname{Pic}^{0}(\widetilde{X}_s) \to 0.
\]
可以证明(Serre，同上)$\operatorname{Pic}^{0}(X_s)$ 等于 $X_s \setminus \{P\}$ 上次数为零的因子群除以形如 $(f)$($f(P) = 1$)的主因子.序列中的第一个映射可以如下描述：设 $a \in \mathbb{G}_{m}(k) = k^{\times}$；因为 $\widetilde{X}_s$ 上的正则函数分离点，存在 $f \in k(\widetilde{X}_s) = k(X_s)$，使得 $f(P_1) = a$ 且 $f(P_2) = 1$；$a$ 的像是 $(f)$.

注意对序列中映射的描述涉及对映到奇异点 $P$ 的点 $P_1$，$P_2$ 排序的选取.相反的选取给出映射的负值.这个符号不确定性贯穿整个理论.

由上述正合序列与 Kummer 序列的上同调序列，得到正合序列
\[
0 \to \mathbb{Q}_{\ell}(1) \to H^{1}(X_s, \mathbb{Q}_{\ell}(1)) \to H^{1}(\widetilde{X}_s, \mathbb{Q}_{\ell}(1)) \to 0,
\]
从而(扭转 $-1$，即与 $\mathbb{Q}(-1)$ 作张量积)正合序列
\[
0 \to \mathbb{Q}_{\ell} \to H^{1}(X_s, \mathbb{Q}_{\ell}) \to H^{1}(\widetilde{X}_s, \mathbb{Q}_{\ell}) \to 0.
\]
特别地，$H^{1}(X_s, \mathbb{Q}_{\ell})$ 的维数为 $2g - 1$.写 $E_s$ 为 $H^{1}(X_s, \mathbb{Q}_{\ell}) \to H^{1}(\widetilde{X}_s, \mathbb{Q}_{\ell})$ 的核.它是消失循环群.\footnote{拓扑学家有一种可视化方法，其中(同调中的)消失循环在族中移动并在某些点消失.我一直无法理解这幅图，但期待看到这部电影.}注意 $E_s \cong \mathbb{Q}_{\ell}$(同构良定义到符号)---我们用 $\delta_s$ 表示 $E_s$ 中对应 1 的元素.

设 $V = R^{1}\varphi_{*}\mathbb{Q}_{\ell}$.则对所有 $t \notin S$，$V_{\bar{\xi}} \cong H^{1}(X_t, \mathbb{Q}_{\ell})$，且 $V_s \cong H^{1}(X_s, \mathbb{Q}_{\ell})$.设 $V = V_{\bar{\xi}}$.可以证明逆特殊化映射 $V_s \to V_{\bar{\xi}}$ 单射，像为 $V^{I_s}$，其中 $I_s \subset \pi_1(U)$ 是 $s$ 处的惯性群.此外，在杯积配对
\[
H^{1}(X_t, \mathbb{Q}_{\ell}) \times H^{1}(X_t, \mathbb{Q}_{\ell}) \to H^{2}(X_t, \mathbb{Q}_{\ell}),
\]
\[
\parallel \qquad \parallel \qquad \parallel
\]
\[
V \qquad V \qquad \mathbb{Q}_{\ell}(-1)
\]
中，$E_s$ 是 $V_s$ 的精确零化子.换言之，序列
\[
0 \to H^{1}(X_s, \mathbb{Q}_{\ell}) \to H^{1}(X_{\bar{\xi}}, \mathbb{Q}_{\ell}) \xrightarrow{x \mapsto x \smile \delta} \mathbb{Q}_{\ell}(-1) \to 0
\]
正合.Lefschetz 铅笔的理论表明对纤维维数为奇数的任何 Lefschetz 铅笔有类似序列.

\subsection*{Lefschetz 铅笔中的上同调}

设 $\varphi \colon X^{*} \to \mathbb{P}^{1}$ 是由 Lefschetz 铅笔产生的映射，$S \subset \mathbb{P}^{1}$ 是使 $X_s$ 奇异的 $\mathbb{P}^{1}$ 的子集.设 $n = 2m + 1$ 是 $\varphi$ 的纤维的维数---于是我们假定纤维维数为奇数.我们设
\[
U = \mathbb{P}^{1} \setminus S;
\]
\[
\pi_1(U) = \pi^{\mathrm{tame}}_1(U, \bar{\xi});
\]
\[
I_s = s \text{ 处的驯基本群}
\]
($\pi_1(U)$ 的子群，良定义到共轭).

设 $V = (R^{n}\varphi_{*}\mathbb{Q}_{\ell})_{\bar{\xi}}$，$V(r) = V \otimes \mathbb{Q}_{\ell}(r)$.则下面的事实成立.
\begin{itemize}
  \item 对 $r \neq n, n+1$，层 $R^{r}\varphi_{*}\mathbb{Q}_{\ell}$ 在 $\mathbb{P}^{1}$ 上局部常值(从而常值)，即它们``看不到''奇点.
  \item 紧合-光滑基变换定理蕴含 $R^{n}\varphi_{*}\mathbb{Q}_{\ell}|_{U}$ 局部常值，因此 $V$ 是 $\pi_1(U, \bar{\xi})$-模.事实上，作用穿过驯基本群 $\pi_1(U)$.
  \item 对每个 $s \in S$，有``消失循环'' $\delta_s \in V(m)$(良定义到符号).设 $E(m)$ 是 $V(m)$ 中由诸 $\delta_s$ 生成的子空间---则 $E \subset V$ 是消失循环空间.
\end{itemize}
若一个消失循环为零，则它们全为零.下面我假定没有消失循环为零.这是典型情形，另一种情形中 Riemann 猜想的证明容易得多.在这个假设下，$R^{n+1}\varphi_{*}\mathbb{Q}_{\ell}$ 也常值.唯一还需要理解的层是 $R^{n}\varphi_{*}\mathbb{Q}_{\ell}$.

\begin{itemize}
  \item 对每个 $s \in S$，序列
  \[
  0 \to H^{n}(X_s, \mathbb{Q}_{\ell}) \to H^{n}(X_{\bar{\xi}}, \mathbb{Q}_{\ell}) \xrightarrow{x \mapsto (x \smile \delta)} \mathbb{Q}_{\ell}(m - n) \to 0
  \]
  正合.这里 $(x \smile \delta)$ 是 $(x, \delta)$ 在配对
  \[
  H^{n}(X_{\bar{\xi}}, \mathbb{Q}_{\ell}) \times H^{n}(X_{\bar{\xi}}, \mathbb{Q}_{\ell}(m)) \to H^{2n}(X_{\bar{\xi}}, \mathbb{Q}_{\ell}(m)) \cong \mathbb{Q}_{\ell}(m - n)
  \]
  下的像.
  \item $\gamma_s \in I_s$ 的元素按如下规则作用于 $V$：
  \[
  \gamma_s(x) = x \pm t(\gamma_s)(x \smile \delta_s)\delta_s
  \]
  (Picard-Lefschetz 公式).这里 $t$ 是前面定义的映射 $I_s \to \mathbb{Z}_{\ell}(1)$，因此 $t(\gamma_s)(x \smile \delta_s)\delta_s \in V(1 + (m - n) + m) = V$.符号只依赖于 $n$ 模 4.
\end{itemize}
由这两个论断推出，对所有 $s \in S$，映射 $H^{n}(X_s, \mathbb{Q}_{\ell}) \to H^{n}(X_{\bar{\xi}}, \mathbb{Q}_{\ell})$ 单射，像为 $H^{n}(X_{\bar{\xi}}, \mathbb{Q}_{\ell})^{I_s}$，从而 $R^{n}\varphi_{*}\mathbb{Q}_{\ell} \xrightarrow{\sim} j_{*}j^{*}R^{n}\varphi_{*}\mathbb{Q}_{\ell}$(见命题~\ref{prop:32.1}).

设 $\psi$ 是由杯积定义的形式 $V \times V \to \mathbb{Q}_{\ell}(-n)$.它非退化(Poincar\'e 对偶)且斜对称(因为 $n$ 奇数).

\begin{proposition}\label{prop:32.2}
消失循环空间 $E$ 在 $\pi_1$ 作用下稳定.设 $E^{\perp}$ 是 $E$ 在配对 $\psi$ 下在 $H^{n}(X_{\bar{\xi}}, \mathbb{Q}_{\ell})$ 中的正交补.则 $E^{\perp} = H^{n}(X_{\bar{\xi}}, \mathbb{Q}_{\ell})^{\pi_1}$.
\end{proposition}

\begin{proof}
两个论断都由 Picard-Lefschetz 公式推出.例如，把 Picard-Lefschetz 公式中的 $x$ 取为 $\delta_{s'}$，
\[
\gamma_s(\delta_{s'}) = \delta_{s'} \pm t(\gamma_s)(\delta_{s'} \smile \delta_s)\delta_s, \qquad \gamma_s \in I_s,
\]
蕴含第一个论断，而
\[
\gamma_s(x) - x = \pm t(\gamma_s)(x \smile \delta_s)\delta_s
\]
蕴含第二个论断(因为 $I_s$ 生成 $\pi_1$ 且诸 $\delta_s$ 非零).
\end{proof}

\begin{theorem}\label{thm:32.3}[经典情形中的 Lefschetz]
消失循环在 $\pi_1$ 作用下共轭到符号，即给定 $s, s' \in S$，存在 $\gamma \in \pi_1$，使得 $\gamma\delta_{s'} = \delta_s$ 或 $-\delta_s$.
\end{theorem}

\begin{proof}
(梗概) 设 $D$($= \mathbb{P}^{1}$)是 Lefschetz 铅笔构造中的直线.回忆 $D \subset \check{\mathbb{P}}^{N}$ 且 $S = D \cap \check{X}$.

事实 1：映射 $\mathbb{P}^{1} \setminus S \to \check{\mathbb{P}}^{m} \setminus \check{X}$ 诱导满射 $\pi_1(\mathbb{P}^{1} \setminus S) \to \pi_1(\check{\mathbb{P}}^{N} \setminus \check{X})$.

这由 Bertini 定理推出，即若 $Y$ 不可约且 $Y \to \check{\mathbb{P}}$ 支配，则 $Y$ 到 $\check{\mathbb{P}}$ 中``一般''直线的拉回也不可约.

事实 2：$\pi_1(\mathbb{P}^{1} \setminus S)$ 在 $V$ 上的作用穿过 $\pi_1(\check{\mathbb{P}}^{N} \setminus \check{X})$.

事实 3：选取 $\mathbb{Z}_{\ell}(1)$ 的生成元 $\delta$，对每个 $s \in S$ 选取 $\gamma_s \in I_s$，使得 $t(\gamma_s) = \delta$.则诸 $\gamma_s$ 在 $\pi_1(\check{\mathbb{P}}^{N} \setminus \check{X})$ 中共轭.

在经典情形中，事实 3 有简单的几何证明.抽象情形更困难.

现在我们用 Picard-Lefschetz 公式完成证明.首先注意公式
\[
\gamma_s x - x = \pm t(\gamma_s)(x \smile \delta_s)\delta_s, \qquad x \in V,
\]
把 $\delta_s$ 确定到符号.其次注意对 $\gamma \in \pi_1(\mathbb{P}^{1} \setminus S)$，
\[
(\gamma\gamma_s\gamma^{-1})x = \gamma(\gamma^{-1}x \pm t(\gamma_s)(\gamma^{-1}x \smile \delta_s)\delta_s) = x \pm t(\gamma_s)(x \smile \gamma\delta_s)\gamma\delta_s.
\]
事实上，这个公式对 $\gamma \in \pi_1(\check{\mathbb{P}}^{N} \setminus \check{X})$ 成立，因此``事实''完成证明.
\end{proof}

\begin{corollary}\label{cor:32.4}
空间 $E/E \cap E^{\perp}$ 是绝对简单的 $\pi_1(\mathbb{P}^{1} \setminus S)$-模，即即使在与 $\mathbb{Q}_{\ell}$ 的扩张域作张量积后，它也不含被 $\pi_1(\mathbb{P}^{1} \setminus S)$ 稳定的非零真子模.
\end{corollary}

\begin{proof}
设 $x \in E$.若 $x \notin E^{\perp}$，则存在 $s$，使得 $x \smile \delta_s \neq 0$，从而公式 $\gamma_s x = x \pm t(\gamma_s)(x \smile \delta_s)\delta_s$ 表明由 $x$ 与它在 $\pi_1(U)$ 的元素下的所有变换张成的空间含 $\delta_s$，从而含所有消失循环.因此它等于 $E$.这个论证对 $\mathbb{Q}_{\ell}$ 的任何扩张域都成立.
\end{proof}

\begin{remark}\label{rem:32.5}
设 $\psi$ 是由杯积定义的形式 $V \times V \to \mathbb{Q}_{\ell}(-n)$；它斜对称(因为 $n$ 奇数)，且在 $E/E \cap E^{\perp}$ 上非退化.它被单值群 $\pi_1(U_0)$ 保持，因此 $\pi_1(U)$ 映到 $\operatorname{Sp}(E/E \cap E^{\perp}, \psi)$ 中.
\end{remark}

\begin{theorem}\label{thm:32.6}[Kazhdan 与 Margulis]($n$ 奇数.) $\pi_1(\mathbb{P}^{1} \setminus S, \bar{\xi})$ 在 $\operatorname{Sp}(E/E \cap E^{\perp}, \psi)$ 中的像开.
\end{theorem}

这由上面综述的结果与下一个引理推出.注意，因为 $\pi_1(\mathbb{P}^{1} \setminus S, \bar{\xi})$ 紧，它在 $\operatorname{Sp}(E/E \cap E^{\perp}, \psi)$ 中的像闭.

\begin{lemma}\label{lem:32.7}
设 $\psi$ 是 $\mathbb{Q}_{\ell}$ 上向量空间 $W$ 上的非退化形式.设 $G \subset \operatorname{Sp}(W, \psi)$ 是满足如下条件的闭子群：
\begin{enumerate}
  \item $W$ 是简单 $G$-模；
  \item $G$ 由形如
  \[
  x \mapsto x \pm \psi(x, \delta)\delta
  \]
  ($\delta \in G$ 适当)的自同构拓扑生成.
\end{enumerate}
则 $G$ 含 $\operatorname{Sp}(W, \psi)$ 的开子群.
\end{lemma}

\begin{proof}
将需要一点 $\mathbb{Q}_{\ell}$ 上 Lie 群的理论，对此我参考 Serre, J.-P., \emph{Lie Algebras and Lie Groups}, Benjamin, 1965.因为 $G$ 在 $\operatorname{Sp}$ 中闭，它是 $\mathbb{Q}_{\ell}$ 上的 Lie 群(由 Cartan 定理的 $\ell$ 进类比).设 $L = \mathrm{Lie}(G)$(等于 $G$ 在 1 处的切空间).为证明引理，只需证明 $L$ 等于 $\mathrm{Lie}(\operatorname{Sp})$，因为指数映射把 Lie 群的 Lie 代数中 0 的任意足够小邻域映到 Lie 群中 1 的邻域上.

还有映射 $\log \colon G \to \mathrm{Lie}\, G$(定义在 1 的邻域上).设 $\delta \in W$，$\alpha$ 是自同态 $x \mapsto x \pm \psi(x, \delta)\delta$.则
\[
\log(\alpha) = \log(1 - (1 - \alpha)) = -\sum \frac{(1 - \alpha)^{n}}{n}.
\]
但 $(1 - \alpha)(x) = \pm \psi(x, \delta)\delta$，因此因为 $\psi(\delta, \delta) = 0$, $(1 - \alpha)^{2} = 0$.因此 $\mathrm{Lie}(G)$ 含自同态
\[
N(\delta) \colon x \mapsto \pm \psi(x, \delta)\delta,
\]
且由它们生成.于是，如下关于 Lie 代数的论断将完成证明.
\end{proof}

\begin{lemma}\label{lem:32.8}
设 $W$ 是特征零域 $k$ 上的向量空间，$\psi$ 是 $W$ 上的非退化形式.设 $L$ 是 $\operatorname{Sp}(W, \psi)$ 的 Lie 代数的子-Lie 代数，使得
\begin{enumerate}
  \item $W$ 是简单 $L$-模；
  \item $L$ 由某些形如 $N(\delta)$ 的自同态($\delta \in W$ 适当)生成.
\end{enumerate}
则 $L$ 等于 $\operatorname{Sp}(W, \psi)$ 的 Lie 代数.
\end{lemma}

证明省略(暂时)---它大约两页，且是初等的.
